% Generated draft: requires mathematical and source review.
\documentclass[11pt]{article}
\usepackage[margin=25mm]{geometry}
\usepackage{fontspec}
\setmainfont{Latin Modern Roman}
\usepackage{amsmath,amssymb,mathtools,unicode-math}
\setmathfont{Latin Modern Math}
\usepackage{enumitem,needspace,xurl,hyperref}
\setlength{\parindent}{0pt}
\setlength{\parskip}{5pt}
\setlength{\emergencystretch}{3em}
\newcommand{\N}{\mathbb{N}}
\newcommand{\Z}{\mathbb{Z}}
\newcommand{\Q}{\mathbb{Q}}
\newcommand{\R}{\mathbb{R}}
\newcommand{\C}{\mathbb{C}}
\newcommand{\F}{\mathbb{F}}
\newcommand{\A}{\mathbb{A}}
\newcommand{\PP}{\mathbb{P}}
\newcommand{\E}{\mathbb{E}}
\newcommand{\eps}{\varepsilon}
\newcommand{\dd}{\,\mathrm{d}}
\begin{document}
\section*{UnsolvedMath ID: 2303029}
\textbf{Source problem code:} catalogue entry
\subsection{Definitions and mathematical statement}
For a bounded full-dimensional body $K\subset\R^n$ with uniform unit density, its Newtonian potential, up to an irrelevant nonzero normalization factor, is
\[
U_K(x)=\int_K\|x-y\|^{2-n}\,dy\quad(n\geq3).
\]
In dimension $2$ the corresponding potential uses the logarithmic kernel, for example $U_K(x)=-\int_K\log\|x-y\|\,dy$. A quadratic interior potential means that there is one polynomial $Q\in\R[x_1,\ldots,x_n]$ of degree at most $2$ with $U_K(x)=Q(x)$ throughout the interior of $K$ (with boundary equality interpreted according to the body convention). It is not merely a different quadratic formula near each point.

The source states that ellipsoids have this property and that Nikliborc and Dive independently proved the converse in dimensions $2$ and $3$. H. S. Shapiro asks for the converse in every integer dimension $n>3$:
\[
U_K\text{ is quadratic inside }K
\quad\Longrightarrow\quad
K\text{ is an ellipsoid}.
\]
An ellipsoid is an affine image of a Euclidean ball, equivalently a set of the form $\{x:(x-c)^{\mathsf T}A(x-c)\leq1\}$ with $A$ symmetric positive definite, or its interior if open bodies are used. The source prefers a new method because the earlier proofs involve highly nontrivial calculations; this is a methodological preference rather than an extra hypothesis on $K$.

The supplied word ``body'' does not specify bounded-domain, connectedness, boundary-regularity, or regular-closure assumptions. These matter to exact set equality: the potential only depends on the body up to Lebesgue-null changes, so a conclusion as an exact set cannot be posed for completely arbitrary measurable sets without a convention. No missing body hypotheses are silently supplied.

The research summary claims a solution in all dimensions by Shapiro and later results, describing a ``rat-trap''/harmonic-extension theorem and mentioning Reznick and algebraic approaches. It supplies no precise theorem or reference establishing the required scope. This claimed resolution conflicts with the registry's open status and OPEN-TRIAGE classification; both are preserved as supplied assessments rather than reconciled by an invented literature verification.
\subsection{Short English statement}
The question asks whether a uniform body whose interior Newtonian potential is quadratic must be an ellipsoid in dimensions above $3$, preferably with a less computational proof. The supplied summary claims an all-dimensional solution, but gives no precise citation and conflicts with the open registry label. The body's regularity convention also needs clarification.
\subsection{Sources}
\begin{itemize}
\item Local UnsolvedMath registry, permanent ID 2303029, exact source code catalogue entry, ``Research Problems in Function Theory --- Problem 3.29''; registry status open and research classification OPEN-TRIAGE, despite the supplied summary's SOLVED-IN-LITERATURE assertion.
\item Hayman, \emph{Research Problems in Function Theory} (2018), Problem 3.29, attributed to H. S. Shapiro, \url{https://arxiv.org/abs/1809.07200}. The record identifies extraction from public arXiv source TeX, accessed July 29, 2026; the 2018 update was retained for status triage.
\item Nikliborc and Dive, credited independently for the dimension-$2$ and dimension-$3$ converse. Their parenthetical references are incomplete in the supplied statement; no titles or dates are available here.
\item The supplied research summary attributes an all-dimensional converse to H. S. Shapiro and later work, mentions a ``rat-trap''/harmonic-extension theorem, Reznick, and an algebraic approach to the quadratic-potential problem. It provides no exact publication, theorem number, or hypotheses for these claimed resolutions, so their attribution and scope require verification.
\end{itemize}
\end{document}
